\section*{\S\ 7. Entwicklung der allgemeinen Formen nach Normalformen.}
Wir werden im folgenden (\S\ 8) eine Haupteigenschaft des Faltungsprozesses (41.) ableiten, seine Zusammensetzbarkeit aus Grundfaltungen, bed\"urfen dazu aber der \"Ubertragung eines wichtigen, von Herrn Mertens\footnote{F. Mertens, \emph{\"Uber invariante Gebilde quatern\"arer Formen} (siehe Einleitung), zitiert als ,,Mertens``.} im quatern\"aren Gebiete gegebenen Satzes auf das $n$-n\"are Gebiet. Dieser Satz sagt aus, da\ss\ sich ein beliebiges endliches System von Formen ersetzen l\"a\ss t durch ein \"aquivalentes System von Normalformen. Unter ,,Normalform`` verstehen wir dabei -- in Verallgemeinerung der bin\"aren und tern\"aren Definition\footnote{Study, a. a. O. S. 54.} -- eine Form, deren s\"amtliche in den Koeffizienten lineare invariante Bildungen identisch verschwinden mit Ausnahme der Grundform selbst. Nach Satz VI (\S\ 6) gen\"ugt also eine Normalform $\varphi$ dem System von Differentialgleichungen:
\begin{equation}
 \left(q_{n-\sigma-\lambda}\frac{\partial}{\partial p_\sigma}\middle|\frac{\partial}{\partial q_{n-\tau}}p_{\tau-\lambda}\right)\varphi=0,
 \qquad(\sigma\geqq\tau,\ \lambda=1,2,\ldots,\tau,n-\sigma),
\tag{46}
\end{equation}
wo die Reihen $q_{n-\sigma-\lambda},p_{\tau-\lambda}$ als willk\"urliche Gr\"o\ss en zu betrachten sind. Diese Differentialgleichungen stimmen im quatern\"aren Gebiet -- unter Ber\"ucksichtigung von (39.) f\"ur $\sigma=\tau=2$ -- mit den von Herrn Mertens ohne Einf\"uhrung der willk\"urlichen Reihen $p,q$ gegebenen 10 Differentialgleichungen (a. a. O., Formel 28)) vollst\"andig \"uberein; ihre Kenntnis, zusammen mit Satz VI, erlaubt uns, das dortige Beweisverfahren fast w\"ortlich auf $n$ Variable zu \"ubertragen. Hinzuzuf\"ugen ist einzig der mittels der Zerlegungsidentit\"at (30.) erbrachte Nachweis, da\ss\ die konstruierten Formen Normalformen sind; im quatern\"aren Gebiet ergibt sich dieser Nachweis noch ohne weitere Umformung. Es gen\"ugt deshalb auch eine kurze Darlegung des Beweises, und zwar der Einfachheit halber in dem f\"ur das folgende nur in Betracht kommenden Fall einer einzigen Grundform $f$ und ihrer Formenreihe (vgl. \S\ 6 Schlu\ss), da f\"ur invariante Bildungen beliebigen Grades in den Koeffizienten beliebig vieler Grundformen der Beweis derselbe bleibt.

Der Grundgedanke des Beweises ist, durch Einf\"uhrung der transformierten Koeffizienten eine Form mit $n-1$ Reihen $p_\rho$ zu ersetzen durch Formen mit $n$ den $x$ kogredienten Reihen $\xi$, den Reihen der Transformationsgr\"o\ss en. Aus diesen Formen leiten sich durch invariante Differentiationsprozesse direkt die gesuchten Normalformen ab. Die Entwicklung der allgemeinen Formen nach Normalformen wird vermittelt durch eine Reihenentwicklung f\"ur Formen mit $\rho$ ($\rho\leqq n$) kogredienten Reihen $\xi$; invariante Differentiationsprozesse f\"uhren zu den Formen in den urspr\"unglichen Reihen $p_\rho$ zur\"uck.
